\documentclass[10pt,a4paper]{amsart}
\usepackage[margin=19mm]{geometry}
\usepackage{amsmath,amssymb,mathtools}
\usepackage[scr=rsfs,cal=euler]{mathalfa}
\usepackage{xfrac}
\usepackage[unicode,hidelinks,bookmarks=false]{hyperref}
\newcommand{\ie}{i.e.\ }
\newcommand{\cf}{cf.\ }
\newcommand{\resp}{respectively\ }
\newcommand{\Subsection}{\subsection}
\providecommand{\nobreakdash}{\nobreak\hbox{-}}
\setlength{\emergencystretch}{2em}
\multlinegap0pt
\usepackage{amscd}
\usepackage{amsthm}
\newtheorem{defn}{Definition}
\newtheorem{fact}{Fact}
\newtheorem{thm}{Theorem}
\newtheorem{prop}{Proposition}
\newtheorem{lem}{Lemma}
\newtheorem{conj}{Conjecture}
\newtheorem{slem}{Sublemma}
\newtheorem{cor}{Corollary}
\newtheorem{ass}{Assumption}
\newtheorem{ques}{Question}
\newtheorem*{mlem*}{Main Lemma}

\newcounter{constant} 
\newcommand{\newconstant}[2]{\refstepcounter{constant}\label{#1}}
\newcommand{\useconstant}[2]{C_{\ref{#1}}}
\title[Random walks and the symplectic representation of the braid group]{Random walks and the symplectic representation of the braid group: transience of a polynomial zero set on the symplectic representation}
\author{Marc Soret and Marina Ville}
\date{Source: arXiv:2203.00984v2 (CC BY 4.0)}
\begin{document}
\maketitle
\noindent\emph{Source and licence.} Random walks and the symplectic representation of the braid group, arXiv:2203.00984v2, distributed under the Creative Commons Attribution 4.0 International licence, http://creativecommons.org/licenses/by/4.0/, as recorded in the arXiv OAI metadata for this exact version and shown on the abstract page https://arxiv.org/abs/2203.00984v2. Full rights notice: licenses/arxiv-2203.00984-CC-BY-4.0.txt.\par\smallskip
\noindent\textit{Editorial scope.} Counted unit: the paper\textquotesingle s main theorem (source label \textbackslash{}label\{prop sur les transients\}, source0.tex lines 109--115, printed as Theorem 3 of the article), the transience of the zero set of an integer polynomial evaluated on the symplectic representation $\rho_n$, delivered with its complete original proof at source lines 166--250 as the single primary proof body. The source\textquotesingle s own wording, numbering and notation are reproduced exactly; only the publisher\textquotesingle s \textbackslash{}documentclass\{article\} and its package list are replaced by the AMS amsart wrapper plus the author\textquotesingle s own theorem environments and macros, so that the counted statement prints with the article\textquotesingle s own number (Theorem 3: the two earlier theorems of the paper, Przytycki\textquotesingle s and Malyutin\textquotesingle s, are retained in the context body and carry the counter). One editorial counter command immediately before the proof restores the article\textquotesingle s own numbering for the two theorems that the proof states internally, which the article prints as Theorem 6 and Theorem 7: the article\textquotesingle s Theorems 4 and 5 lie outside this unit, and without the command the two internal theorems would print as Theorem 4 and Theorem 5. Every other printed number of the delivered unit (Theorem 1, Theorem 2, Theorem 3, Lemma 1, equations (1)--(14), the section numbers and the bibliography labels) coincides with the article without any editorial intervention. Required context retained uncounted: the whole of section 1 Preliminaries with the definition of the symplectic representation, of a nondegenerate probability, of convolution and of the right random $\mu$-walk (lines 36--104), and Assumption 1 (lines 105--108). The paper\textquotesingle s own bibliography (lines 618--651) is retained as context so that the literal bracket attributions of the source text remain resolvable. The delivered unit contains no \textbackslash{}cite command: this article marks its citations with literal bracket labels such as [Mal], [A\textquotesingle C], [Di], [As], [L-R], [Mi] rather than with \textbackslash{}cite, so the citation map of the package is empty by construction. Omitted under the declared bounded\_source\_comparison fidelity receipt: the single bibliography line 636, \textbackslash{}bibitem[Ma]\{ma\} J. Ma, which shares the key "ma" with the Malyutin entry at line 638 and would therefore raise a multiply-defined label; the J. Ma reference is not cited anywhere in this unit and the Malyutin entry, which is the one the counted context uses, is retained verbatim. No other source line is omitted. Printed slips of the source are kept literal and are not repaired: the misspelling "nondegerate" (line 96), the phrase "the only normal subgroups of $SL(2,\mathbb{Z})$" where $SL(2,\mathbb{F}_3)$ is meant (line 192), the matrix identity $(s_2s^{-1})^4=Id$ written with a missing subscript (line 217), and the stray \textbackslash{}label\{QQ\} that the source uses twice (lines 475 and 668; neither occurrence lies inside this unit). No mathematics is written, repaired or re-derived; every mathematical symbol in this unit is the author\textquotesingle s own native text.
\section{Preliminaries}
\subsection{Representations of the braid group}\label{paragraphe sur les representations}
If $B(n+1)$ is the group of braids with $n+1$ strands, we denote by 
$\mathcal B_t$ its {\it reduced Burau representation} which represents braids as $n\times n$ matrices with values in $\mathbb{Z}[t,t^{-1}]$. The specialization of $\mathcal B_t$ for $t=-1$, denoted $\mathcal B_{-1}$, is called the {\it integral reduced Burau representation}. The representation $\mathcal B_{-1}$ is closely related to the {\it symplectic representation} of $B(n+1)$ which we now recall, using the description of A'Campo ([A'C] \S  1).\\
We consider a surface $X$ such that $H_1(X,\mathbb{Z})$ is generated by $n$ loops; the generators of $B(n+1)$ act on $H_1(X,\mathbb{Z})$ as Dehn twists w.r.t. corresponding generators of $H_1(X,\mathbb{Z})$. This action is conjugate to the integral reduced Burau representation ${\mathcal B}_{-1}$. \\  
This action of $B(n+1)$ preserves the intersection form $I$ on $H_1(X,\mathbb{Z})$. 
\begin{enumerate}
	\item 
	If $n$ is even, $n=2l$, the form $I$ is symplectic nondegenerate on $H_1(X,\mathbb{Z})$.
	\item
	If $n$ is odd, $n=2l+1$, the kernel of $I$ is a line $K$; the map $\mathcal B_{-1}$ restricts to the identity on $K$ and  acts symplectically on the quotient $H_1(X,\mathbb{Z})/K$ 
\end{enumerate}
Thus, in both cases, we get  a {\it symplectic representation}
\begin{equation}\label{representation symplectique0}
\rho_n:B(n+1)\longrightarrow Sp(2l,\mathbb{Z})
\end{equation} 
\subsection{Random walks on groups}\label{paragraph on random walks}
If $\mu_1,\mu_2$ are probability measures on a discrete group $G$ with finite support, their {\it convolution} is defined for $g\in G$ as 
\begin{equation}\label{definition de la convolution}
(\mu_1\star\mu_2)(g)=\sum_{h\in G}\mu_1(gh^{-1})\mu_2(h)
\end{equation}
\begin{defn}\label{non degenerate}
	A  probability $\mu$ on a discrete group $G$ is nondegenerate if $supp(\mu)$ generates $G$ as a semigroup.\end{defn} Malyutin ([Mal]) defines the {\it  right random $\mu$-walk} as the Markov chain which starts at the identity of $G$ and with transition probability $P(g,h)=\mu(gh^{-1})$. If $X$ is a subset of $G$, the probability that this walk hits $X$ at the $k$-th step is $\mu^{\star k}(X)$. 
\subsection{Signature and determinant of a link}\label{definition de la signature}
These are classical invariants of links and knots. 
If $\gamma$ is a link on $\mathbb{S}^3$ which bounds a Seifert surface $\Sigma$ in $\mathbb{S}^3$, we consider the bilinear form called {\it Seifert form} $$\Phi: H_1(\Sigma,\mathbb{Z})\times H_1(\Sigma, \mathbb{Z})\longrightarrow\mathbb{Z}$$
$$(\alpha,\beta)\mapsto lk(\hat{\alpha},\beta)$$ where $lk$ is the linking number and $\hat{\alpha}$ is the link obtained by pushing $\alpha$ in the direction normal to $\Sigma$ in $\mathbb{S}^3$. Given a base for $H_1(\Sigma, \mathbb{Z})$, we define the matrix $V$ of $\Phi$ and consider the bilinear form defined by the matrix $V+^tV$: its determinant is the {\it determinant} of the link and its {\it signature} is the signature of the link. We point out that the literature contains two different definitions of the signature which coincide up to sign.\\
If $\Delta_L$ is the Alexander polynomial of a link, then $det(L)=\Delta_L(-1)$.\\ 
We recall (cf. [Mu]) that for a knot $K$, 
\begin{equation}
|sign(K)|\leq 2g_4(K)
\end{equation} where $g_4$ denotes the topological $4$-genus.  Thus a slice knot has zero signature but the converse is not true and we will see examples of that below.
\subsection{Positive links and quasipositive braids}
A {\it positive link} is a link which has a diagram with all positive crossings.
\begin{thm}([Pr])
	A positive link has strictly negative signature.
\end{thm} 
So it is interesting to look at the signature of quasipositive braids:
\begin{defn} ([Ru])
	A braid $\beta$ is {\it quasipositive} if it is a product of conjugates of positive braid generators
	\begin{equation}\label{expression de la tresse quasipositive}
	\beta=\prod_{i=1}^k\gamma_{i}\sigma_{j_i}\gamma_{i}^{-1}
	\end{equation}
\end{defn}
If we close the  quasipositive $N$-braid (\ref{expression de la tresse quasipositive}) in a link $\hat{\beta}$, Rudolph proved ([Ru]) that
\begin{equation}\label{quasipositif.link}
\chi_4(\hat{\beta})=N-k
\end{equation}
where $\chi_4$ denotes the largest Euler characteristic of a {\it smooth} surface in $\mathbb{B}^4$ bounded by $\hat{\beta}$.\\
Tanaka ([Ta]) constructed examples of quasipositive braids with zero signature and we will give some more here.

\subsection{Transient sets in the braid groups from quasimorphisms}
 We recall that 
a {\it quasimorphism} on a group $G$ is a function $\Phi:G\longrightarrow \mathbb{R}$ such that there exists a positive constant $D_\Phi$ with
\begin{equation}\label{quasimorphisme}
\forall a,b\in G,\ |\Phi(ab)-\Phi(a)-\Phi(b)|\leq D_\Phi
\end{equation}
where $D_\Phi$ is called the {\it defect} of $\Psi$.\\
Quasimorphisms give exemples of transient sets in the braid groups:
\begin{thm}\label{malyutin}
	([Mal]) Let $G$ be a countable group and $\mu$ a nondegerate probability on $G$ (see Definition \ref{non degenerate}). If  a subset $S$ of $G$ has bounded image under an unbounded quasimorphism $\Phi$, then the probability that the right random $\mu$-walk hits $S$ at the $k$-th step tends to zero as $k$ tends to infinity.
\end{thm}
We will see below how Gambodau-Ghys's formula ([G-G]) for the signature turns it into a quasimorphism.  Several other link invariants such as $\chi_4$ can be used to define   quasimorphisms on the braid group ([Br], [B-K]).\\
Other exemples of transient sets on the braid groups have been constructed by Ito ([It]).

\section*{Acknowledgements}
We thank Tahl Nowik and Chaim Even-Zohar for reading an earlier draft of this paper, Misha Brandenbursky for very helpful conversation on braids and quasimorphisms and Florence Lecomte for necessary advice on algebraic geometry over finite fields. 


\section{The results}
\begin{ass}\label{assump}  The random walks that we consider here are right random $\mu$ walks for a nondegenerate probability $\mu$ on $B(n+1)$. \\
For $B(3)$, we assume moreover that $\mu(\sigma_i)\neq 0$, $\mu(\sigma_i^{-1})\neq 0$ for $i=1,2$.
\end{ass}

\begin{thm}\label{prop sur les transients}
	Let $P$ be a polynomial in $(2l)^2$ variables with coefficients in $\mathbb{Z}$. If $M=(m_{ij})\in Sp(2l,\mathbb{Z})$, we let 
	$P(M)=P(m_{11},...,m_{12l},m_{21},...,m_{22l},m_{2l1},...,m_{2l2l}).$\\
	Suppose that $P$ does not vanish identically on $Sp(2l,\mathbb{Z})$. Then the set
	$$\{\beta\in B(n+1): P(\rho_n(\beta))=0\}$$
	is transient for the right random $\mu$-walk.
\end{thm}

\setcounter{thm}{5}
\section{Random walks: proof of Theorem \ref{prop sur les transients}}
Similarly to Rivin in [Ri], we introduce the finite groups $Sp(2l,\mathbb{Z}_p)$'s for a prime number $p\geq 3$.\\
A theorem of A'Campo  ([A'C], Theorem 1 (1)) shows that the representation $B(n)\longrightarrow Sp(2l,\mathbb{F}_p)$  is surjective. So we consider the surjective map
\begin{equation}\label{representation symplectique}
\Pi_p:B(n)\overset{\rho_n}{\longrightarrow} Sp(2l,\mathbb{Z})
\longrightarrow Sp(2l,\mathbb{F}_p)
\end{equation}   
\begin{lem}\label{uniforme sur Psp} 
	Let $p\geq 3$, Consider the image probability $\Pi_p\mu$ of $\mu$ via $\Pi_p$ on $Sp(2l,\mathbb{F}_p)$, i.e. $\Pi_p(X)=\mu(\Pi_p^{-1}(X))$.  If $m$ tends to infinity, $(\Pi_p\mu)^{\star m}$ converges to the equidistributed measure on $Sp(2l,{\mathbb{F}_p})$. 
\end{lem}
\begin{proof}
The lemma follows from the following two theorems. 
\begin{thm}\label{groupe fini}
	(see for exemple [Di])
	Let  $G$ be a finite group and $\mu$ a probability measure on $G$ such that 
	\begin{enumerate}
		\item 
$G$ is generated as a semigroup by 	$supp(\mu)$
\item $supp(\mu)$ is not included in a coset of $G$ by a non trivial normal subgroup $H$ of $G$.
\end{enumerate}
  Then $\mu^{\star m}$ converges to the equidistributed measure on $G$ as $m$ tends to infinity.
\end{thm}
\begin{thm}\label{psl simple} ([As]) If $n> 2$ or $p> 3$, the group $PSp(2n,{\mathbb{F}_p})$ is simple.
	\end{thm} 
If $H$ is a non trivial normal subgroup of $Sp(2n,\mathbb{F}_p)$, 
Theorem \ref{psl simple} tells us that it maps into $\{Id\}$, hence $H$ is included into $\{kId: k\in \mathbb{F}_p\}$. Thus $supp(\mu)$ cannot be a  coset $gH$ and Theorem \ref{groupe fini} applies.\\
We now show that if $n=2$ and $p=3$ we can also apply Theorem \ref{groupe fini}. We recall that the only non trivial normal subgroup of $PSL(2,\mathbb{F}_3)$ is the Klein group  so the only normal subgroups of $SL(2,\mathbb{Z})$ are 
\begin{itemize}
	\item the center $\{\pm Id\}$
	\item 
	$H=\{\pm Id, \pm\begin{pmatrix}
	0 & 1 \\
	-1 & 0 
	\end{pmatrix}, \pm\begin{pmatrix}
	-1 & -1 \\
	-1 & 1 
	\end{pmatrix}, \pm\begin{pmatrix}
	-1 & 1 \\
	1 & 1 
	\end{pmatrix}\}$
	\end{itemize}
Now for the reduced Burau representation ${\mathcal B}_{-1}:B(3)\longrightarrow SL(2,\mathbb{Z})$, we have
\begin{equation}\label{les generateurs}
s_1={\mathcal B}_{-1}(\sigma_1)=\begin{pmatrix}
1 & 0 \\
-1 & 1 
\end{pmatrix}\ \ \ \ s_2={\mathcal B}_{-1}(\sigma_2)=\begin{pmatrix}
1 & 1 \\
0 & 1 
\end{pmatrix}
\end{equation}
The support of $\Pi_p\mu$ contains the matrices $s_1$ and $s_2^{-1}$. Suppose there exists $g\in SL(2,\mathbb{F}_3)$ such that $s_1$ and $s_2^{-1}$ belong to $gH$; then $s_2s^{-1}=ghg^{-1}$ for some $h\in H$, hence $(s_2s_1^{-1})^4=Id$, which is not true. Thus the assumptions of Theorem \ref{groupe fini} is verified.
\end{proof}
We now investigate the density of the zero set in $Sp(2l, \mathbb{F}_p)$ of the polynomial  $P$ appearing in Theorem \ref{prop sur les transients}.\\
We write $P$ in terms of all the multi-indices, $P(M)=\sum a_I m^I$, where the $a_I$'s belong to $\mathbb{Z}$. If $p$ is an integer, we denote by $P_p$ its reduction modulo $p$. Since $P$ is not identically zero, we derive an infinite set $\mathcal P$ of prime numbers $p$ such that $P_p$ is not identically zero.  To prove the theorem, we need to estimate
\begin{equation}
\frac{|P_p^{-1}(0)|}{|Sp(2l, \mathbb{F}_p)|}
\end{equation}
We recall (see for example [O'M] or [wi]) that 
\begin{equation}\label{cardinal de Sp}
|Sp(2l,\mathbb{F}_p)|=\prod_{m=1}^l\big[(p^{2m}-1)p^{2m-1}\big]\geq \frac{1}{2^{l}}\prod_{s=1}^{2l}p^{s}
= \dfrac{p^{2l^2+l}}{2^{l}}
\end{equation}
To estimate $|P_p^{-1}(0)|$, we use the following result by Lachaud and Rolland 
\begin{thm}
	([L-R]) Let $K$ be the algebraic closure of $\mathbb{F}_p$. Let $X$ be an algebraic variety of $K^n$ of dimension $m$ which is the zero set of a family of polynomials $(f_1,...,f_r)$. Let $d_i=deg(f_i)$. Then 
	$$|X\cap \mathbb{F}_p^n|\leq d_1...d_rp^m$$
\end{thm}
We view $P_p$ as a polynomial with coefficients in $K$ and let 
\begin{equation}
X=Sp(2l,K)\cap P_p^{-1}(0)\subset \mathbb{F}_p^{4l^2}
\end{equation} We know that for a field $K$, $Sp(2l,K)$ is irreducible as an algebraic variety ([Mi]), hence $X$ is of dimension strictly smaller that $Sp(2l,K)$, i.e. $dim(X)\leq 2l^2+l-1$.  Thus, for some constant $C_1(2l,P)$ depending on $2l$ and on the degree of $P$,
\begin{equation}\label{application de lachaud-rolland}
|X\cap \mathbb{F}_p^{4l^2}|\leq C_1(2l,P) p^{2l^2+l-1}
\end{equation} 
Putting (\ref{cardinal de Sp}) and (\ref{application de lachaud-rolland}) together, we derive that  
\begin{equation}\label{estimee de conclusion}
\dfrac{|X|}{|Sp(2l, \mathbb{F}_p)|}\leq \dfrac{C_2(2l,P)}{p}
\end{equation}
for some constant $C_2(2l,P)$. Thus, given a $\epsilon>0$, we pick a prime number $p$ in $\mathcal P$ such that 
\begin{equation}
2\dfrac{C_2(2l,P)}{p}<\epsilon
\end{equation}
There exists an integer $k_0$ such that for every $k>k_0$, we have $$\mu^{\star k}(P^{-1}(0))\leq (\Pi_p\mu)^{\star k}(P_p^{-1}(0))<\epsilon.$$ 
\qed

	 \begin{thebibliography}{ru}
	 	\bibitem[Ac]{ac} J. Achter, {\it The distribution of class groups of functions} Journal of Pure and Applied Algebra 204 (2006) 316 – 333
	 	\bibitem[A-H]{ah} J. Achter, J. Holden {\it Notes on an analogue of the Fontaine-Mazur
	 	conjecture}, Journal de Th\'eorie des Nombres de Bordeaux 15 (2003), 627–637
	 	\bibitem[A'C]{a'c} N. A'Campo, {\it Tresses, monodromie et le groupe symplectique}, Comm. Math. Helvetici (54) 
	 	318-327, (1979)
	 	 \bibitem[As]{as} M. Aschbacher, {\it Finite group theory}, Cambridge Univ. Press, London 1986
	 	 \bibitem[B-K]{bk} M. Brandenbursky, J. Kedra {\it Concordance group and stable commutator length in braid groups}, Algebraic \& Geometric Topology (5), 2861-2886, (2015).
	 	 \bibitem[Br]{br} M. Brandenbursky, {\it On quasi-morphisms from knot and braid invariants}, Jour. of Knot Theory and its Ramifications, vol. 20, No. 10 (2011), 1397-1417.
	 	 \bibitem[Di]{di} P. Diaconis, {\it Representations in probability and statistics}, IMS, Hayward, 1986
	 	 \bibitem[G-G]{rr} J.-M. Gambaudo, E. Ghys {\it Braids and signatures}, Bull. SMF 133 (2005), 541-579
	 	 \bibitem[It]{it} T. Ito, {\it On a structure of random open books and closed braids}, Proc. Japan Acad. 91, Ser. A (2015) 
 	\bibitem[K-T]{kt} C. Kassel, V. Turaev, {\it Braid groups}, Graduate texts in maths, Springer-Verlag New York (2008)
 %	\bibitem[Ko]{ko} P. Kohn, {\it Unlocking two component links}, Osaka Jour. Math. 30 (1993) 741-752
 \bibitem[La]{la} C. Lamm, {\it Symmetric unions and ribbons knots}, Osaka  J. Math. 37  (2000),  537-550
 		\bibitem[L-O]{la0} C.  Lamm, D. Obermeyer {\it Billiard knots in a cylinder} J. Knot Theory and its Ramifications 8(3) (1999) 353-366.
 		\bibitem[L-R]{lar} G. Lachaud, R. Rolland {\it On the number of points of algebraic sets over finite fields}, Journal of Pure and Applied Algebra
 		219 (11) ,  5117-5136 (2015)

 	%	On the number of points of algebraic sets over finite fields
 		\bibitem[Mal]{ma} A. V. Malyutin, {\it Quasimorphisms, random walks, and transient subsets in countable groups}, 	
 		Jour. of Math. Sciences 181(6) 871--885  (2012)
 			\bibitem[Mi]{mi} J.S. Milne, {\it Algebraic groups}, Cambridge Studies in Advanced Mathematics (170), CUP 2017
 		\bibitem[Mu]{mu}	K. Murasugi, {\it On a certain numerical invariant of link types}, Trans. Amer. Math.
 			Soc. 117 (1965), 387-422. 
 			\bibitem[O'M]{om} O. T. O'Meara, {\it Symplectic groups}, AMS, Providence (1978)
 				\bibitem[Pr]{pr} J. H. Przytycki {\it Positive knots have negative signature}, Bull. Ac. Pol.: Math. 37, 1989, 559-562.
 		\bibitem[Ri]{ri2} I. Rivin, {\it Generic phenomena in groups} in {\it Thin groups and superstrong approximation} MSRI Publications, Volume 61 (2013) 
 		\bibitem[Ro]{ro} D. Rolfsen, {\it Knots and links}, AMS Chelsea Publishing (2003) 
 \bibitem[Ru]{ru} L. Rudolph, {\it Quasipositivity as an obstruction to sliceness}, Bull. Amer. Math. Soc. 29 (1993), 51-59.
 	\bibitem[S-V]{rv} M. Soret, M. Ville, {\it Lissajous-toric knots},  Jour. of Knot Theory and Its Ramifications 29(1) (2020) 
 	\bibitem[Ta]{rt} T. Tanaka, {\it Unknotting numbers  of quasipositive knots}, Topology and its applications 88 (1998) 239-246
 	\bibitem[wi]{wi} https://groupprops.subwiki.org/wiki/Order$\_$formulas$\_$for$\_$symplectic$\_$groups
 \end{thebibliography}

\end{document}
